\booksection{บทที่ 1 ภาคผนวก: หัวข้อพิเศษระบบปฏิบัติการฮาร์โมนีโอเอส}

\chapterintro{HarmonyOS เป็นระบบปฏิบัติการที่หัวเว่ยพัฒนาขึ้นเอง ตั้งแต่ปี 2019 ที่เปิดตัวจนถึงปัจจุบัน ได้ผ่านหลายเวอร์ชันใหญ่ มันไม่ใช่แค่ “อีกหนึ่งทดแทนแอนดรอยด์” แต่เป็นสิ่งใหม่ที่มีปรัชญาการออกแบบที่เป็นเอกลักษณ์ มาทำความเข้าใจระบบปฏิบัติการที่กำลังเติบโตนี้กัน}

\booksubsection{HarmonyOS คืออะไร? ทำไมมันจึงแตกต่าง?}

แนวคิดหลักของ HarmonyOS คือ “การเชื่อมต่อสรรพสิ่งเข้าด้วยกัน” ต่างจากระบบปฏิบัติการดั้งเดิมที่ให้บริการอุปกรณ์เดียวเป็นหลัก HarmonyOS ออกแบบให้รองรับการทำงานร่วมกันของหลายอุปกรณ์ตั้งแต่ต้น

\textbf{สถาปัตยกรรมแบบกระจายศูนย์}เป็นวิญญาณของ HarmonyOS พูดง่ายๆ มันรวมหลายอุปกรณ์เข้าด้วยกันเป็น “ซูเปอร์เทอร์มินัล” โทรศัพท์ แท็บเล็ต คอมพิวเตอร์ นาฬิกา ทีวี ของคุณไม่ได้เป็นเอกเทศแยกขาดจากกันอีกต่อไป แต่เป็น “อวัยวะ” ชุดเดียวกันที่เรียกใช้ความสามารถของกันและกันได้

ยกตัวอย่าง: คุณดูวิดีโอบนมือถือแล้วอยากดูบนจอใหญ่ วิธีดั้งเดิมคือแคสต์หน้าจอ แต่บน HarmonyOS คุณส่งวิดีโอ “ไหลต่อ” ไปที่ทีวีได้เลย แล้วโทรศัพท์กลายเป็นรีโมต นี่ไม่ใช่การแคสต์ธรรมดา แต่เป็นการที่แอปทำงานข้ามอุปกรณ์จริงๆ

\booksubsection{HarmonyOS 4.3: เวอร์ชันที่พัฒนาแล้วเต็มที่บนพื้นฐาน AOSP}

HarmonyOS 4.3 เป็นเวอร์ชันที่ถูกติดตั้งใช้งานแพร่หลายบนโทรศัพท์ Huawei ปัจจุบัน ต้องอธิบายให้เข้าใจว่า HarmonyOS 4.3 รองรับโครงสร้างโค้ดของโครงการแอนดรอยด์โอเพนซอร์ส (AOSP) จึงใช้แอปแอนดรอยด์ได้อย่างครบถ้วน — App ส่วนใหญ่ที่คุณดาวน์โหลดจาก AppGallery ของ Huawei แท้จริงยังเป็นแอปแอนดรอยด์ที่รันผ่านชั้นความเข้ากันได้ของ HarmonyOS แต่ HarmonyOS 4.3 ไม่ใช่การคลุมเปลือกแอนดรอยด์แบบง่ายๆ มันเพิ่มความสามารถแบบกระจายศูนย์และเทคโนโลยีเอนจิน Ark ที่ Huawei พัฒนาเองอีกจำนวนมากบนพื้นฐาน AOSP

\subsubsection*{คุณสมบัติหลัก:}

\begin{itemize}
  \item \textbf{ความเข้ากันได้กับแอปแอนดรอยด์}: เพราะพัฒนาบนพื้นฐาน AOSP HarmonyOS 4.3 จึงรันแอปแอนดรอยด์ได้เกือบทั้งหมด ผู้ใช้แทบไม่รู้สึกถึงความต่าง
  \item \textbf{เอนจิน Ark}: คอมไพเลอร์และรันไทม์ที่ Huawei พัฒนาเอง ทำให้แอปทำงานลื่นขึ้น ประหยัดไฟขึ้น เทียบกับเวอร์ชันก่อนหน้า ความลื่นไหลดีขึ้นชัดเจน
  \item \textbf{ซูเปอร์เทอร์มินัล}: ในศูนย์ควบคุมจะเห็นอุปกรณ์ HarmonyOS อื่นๆ รอบตัวได้อย่างชัดเจน ลากเพื่อเชื่อมต่อได้เลย เรียกใช้คีย์บอร์ดเมาส์ของคอมพิวเตอร์จากโทรศัพท์ เรียกกล้องของโทรศัพท์จากแท็บเล็ต สะดวกมาก
  \item \textbf{การ์ดบริการ}: คล้ายวิดเจ็ตเล็กของ iOS แต่ยืดหยุ่นกว่า ปัดขึ้นบนไอคอนแอปเพื่อขยายการ์ด แสดงข้อมูลสำคัญโดยไม่ต้องเปิดแอป
  \item \textbf{โหมดบริสุทธิ์}: อนุญาตให้ติดตั้งแอปจาก AppGallery ทางการเท่านั้น ลดความเสี่ยงจากซอฟต์แวร์ประสงค์ร้าย
  \item \textbf{ศูนย์ความเป็นส่วนตัว}: รวมศูนย์จัดการสิทธิ์ของแอป ดูบันทึกพฤติกรรมการใช้งานของแอป
\end{itemize}

\subsubsection*{จุดเด่นของอินเทอร์เฟซผู้ใช้:}

\begin{itemize}
  \item \textbf{ล็อกสกรีนสไตล์นิตยสาร}: หน้าล็อกแสดงภาพสวยและบทความ ปัดออกก็เข้าสู่หน้าจอหลัก
  \item \textbf{การ์ดอเนกประสงค์}: ไอคอนบนหน้าจอหลักขยายเป็นการ์ดได้ แสดงข้อมูลได้มากขึ้น
  \item \textbf{สถานีรับส่งกลาง}: พื้นที่จัดเก็บชั่วคราวแห่งหนึ่ง เอาข้อความ รูปภาพ ไฟล์มาพักไว้ได้ แล้วลากไปใช้ในแอปไหนก็ได้ การแชร์ข้ามแอปจึงลื่นสุดๆ
  \item \textbf{การค้นหาอัจฉริยะ}: ค้นหาทั้งระบบ ค้นแอป ไฟล์ การตั้งค่า และหน้าเว็บได้
\end{itemize}

\booksubsection{สาม HarmonyOS NEXT (HarmonyOS 6): จุดเริ่มต้นที่แท้จริงของระบบปฏิบัติการที่ Huawei พัฒนาเอง}

HarmonyOS NEXT (HarmonyOS NEXT 5.0) กับ HarmonyOS 4.3 มีความแตกต่างกันโดยรากฐาน — มันใช้ไมโครเคอร์เนล HarmonyOS ที่ Huawei พัฒนาเอง ไม่พึ่งพาเคอร์เนล Linux และโค้ด AOSP อีกต่อไป และ\textbf{ไม่รองรับแอปแอนดรอยด์โดยตรงอีกต่อไป}นี่คือก้าวสำคัญที่ทำให้ HarmonyOS ผ่านจาก “รูปแบบเปลี่ยนผ่านที่รองรับแอนดรอยด์” ไปสู่ “ระบบปฏิบัติการที่เป็นอิสระโดยสมบูรณ์” อย่างเป็นทางการ

\subsubsection*{การเปลี่ยนแปลงครั้งใหญ่:}

\begin{itemize}
  \item \textbf{เคอร์เนลที่พัฒนาเองโดยสมบูรณ์}: HarmonyOS NEXT ใช้เคอร์เนลที่ Huawei พัฒนาเอง ไม่มีโค้ด AOSP ตกค้างอีกเลย ระบบจึงเรียบง่ายขึ้น ปลอดภัยขึ้น มีประสิทธิภาพมากขึ้น
  \item \textbf{ไม่รองรับแอปแอนดรอยด์อีกต่อไป}: นี่คือการเปลี่ยนแปลงที่ใหญ่ที่สุด HarmonyOS NEXT รันได้แค่แอป HarmonyOS ดั้งเดิม ติดตั้งและรันไฟล์ APK ไม่ได้ ก่อนอัปเกรดโปรดยืนยันว่าแอปที่คุณใช้บ่อยมีเวอร์ชันดั้งเดิมบน HarmonyOS แล้ว
  \item \textbf{ระบบนิเวศแอปดั้งเดิม}: Huawei กำลังผลักดันให้นักพัฒนาย้ายแอปมาสู่แพลตฟอร์ม HarmonyOS ดั้งเดิมอย่างเต็มกำลัง ตอนนี้แอปหลักอย่าง WeChat Alipay Taobao ฯลฯ เปิดตัวเวอร์ชันดั้งเดิมบน HarmonyOS แล้ว แต่แอปจำเพาะบางตัวอาจยังปรับตัวไม่เสร็จ
  \item \textbf{สถาปัตยกรรมความปลอดภัย Star Shield}: ระบบความปลอดภัยที่ออกแบบใหม่ทั้งหมด ปกป้องครบทั้งสายตั้งแต่ฮาร์ดแวร์ถึงซอฟต์แวร์ รองรับการควบคุมสิทธิ์ที่ละเอียดกว่าเดิม เช่น "การระบุตำแหน่งแบบเบลอ" — แอปจะรู้แค่ว่าคุณอยู่แถวไหนโดยประมาณ ไม่รู้ตำแหน่งที่แม่นยำ
  \item \textbf{AI หลอมรวมเชิงลึก}: ความสามารถ AI ระดับระบบ ไม่ใช่แค่มี AI ในแอปใดแอปหนึ่ง เช่น AI คำบรรยายแปลวิดีโอได้แบบเรียลไทม์ AI ตัดรูปแยกตัวหลักได้อย่างฉลาด
  \item \textbf{ไหลต่อข้ามอุปกรณ์ไร้รอยต่อ}: ความสามารถแบบกระจายศูนย์ที่ล้ำหน้าขึ้นไปอีก สถานะแอปสลับข้ามอุปกรณ์ได้ไร้รอยต่อ เอกสารที่แก้อยู่บนโทรศัพท์ พอเดินไปหน้าคอมพิวเตอร์ก็เปิดต่อบนคอมพิวเตอร์ให้เอง
  \item \textbf{Meta Service (บริการแบบทันที)}: บริการเบาๆ ที่ใช้ได้โดยไม่ต้องติดตั้ง เช่น สแกนคิวอาร์โค้ดแล้วเปิดฟังก์ชันนั้นๆ ได้ทันที ไม่ต้องดาวน์โหลดแอปก่อน
\end{itemize}

\begin{tcbwarn}{ต้องอ่านก่อนอัปเกรดเป็น HarmonyOS NEXT}
หลังอัปเกรดเป็น HarmonyOS NEXT โทรศัพท์ของคุณจะรันแอปแอนดรอยด์ไม่ได้อีก โปรดยืนยันก่อนอัปเกรดว่า: (1) แอปที่ใช้บ่อยมีเวอร์ชันดั้งเดิมบน HarmonyOS แล้ว (2) งานหรือชีวิตของคุณไม่มีสถานการณ์ที่จำเป็นต้องพึ่งแอปแอนดรอยด์ และหลังอัปเกรดจะย้อนกลับไป HarmonyOS 4.3 ไม่ได้

\end{tcbwarn}

\subsubsection*{การเปลี่ยนแปลงตรรกะการใช้งาน:}

\begin{itemize}
  \item \textbf{การนำทางด้วยท่าทาง}: ปัดจากขอบซ้ายขวาเพื่อย้อนกลับ ปัดขึ้นจากด้านล่างเพื่อกลับหน้าจอหลัก ปัดขึ้นแล้วค้างไว้เพื่อเข้าหน้าจอหลายงาน
  \item \textbf{หลายหน้าต่างอัจฉริยะ}: รองรับหลายรูปแบบ ทั้งแบ่งจอ หน้าต่างลอย และมุมมองคู่ขนาน ฯลฯ
  \item \textbf{การโทรผ่าน MeeTime}: บริการวิดีโอคอลของ Huawei รองรับการแชร์หน้าจอ ปรับหน้าให้สวย และเบลอพื้นหลัง
\end{itemize}

\booksubsection{ระบบนิเวศ HarmonyOS: อุปกรณ์และแอปพลิเคชัน}

ฮาร์โมนีโอเอสไม่ตัวใช่โทรศัพท์มือถือระบบ, แต่ใช่1ชิ้นคลุมฝามากชนิดอุปกรณ์ของระบบนิเวศ

\subsubsection*{รองรับของอุปกรณ์ประเภท:}

\begin{itemize}
  \item \textbf{โทรศัพท์มือถือและราบบอร์ด}:แก่นกลางอุปกรณ์, การทำงานที่สุดเสร็จทั้งหมด
  \item \textbf{สมาร์ตวอทช์และสมาร์ตแบนด์}:โดดเดี่ยวตั้งทำงานหรือและโทรศัพท์มือถือประสานเดียวกัน
  \item \textbf{จออัจฉริยะ (ทีวี)}:บ้านครอบครัวศูนย์ความบันเทิง
  \item \textbf{อัจฉริยะลำโพง}:เสียงพูดส่งซึ่งกันและกันเข้าปาก
  \item \textbf{รถยนต์}: ห้องโดยสาร HarmonyOS สัมผัสประสบการณ์ห้องขับขี่อัจฉริยะ
  \item \textbf{สมาร์ทโฮม}: ตู้เย็น แอร์ เครื่องซักผ้า ฯลฯ
\end{itemize}

\subsubsection*{สถานะปัจจุบันของระบบนิเวศแอป:}

AppGallery ของ HarmonyOS มีแอปให้เลือกใช้งานนับล้าน โดยแอปหลักๆ จะมีให้ครบโดยพื้นฐานแล้ว Huawei ยังผลักดันให้นักพัฒนาสร้างแอป HarmonyOS เดิมอย่างจริงจัง โดยถึงกลางปี 2026 แอปชั้นนำอย่าง WeChat, Alipay, Taobao, Douyin, Meituan, JD ฯลฯ ต่างก็ได้เปิดตัวเวอร์ชัน HarmonyOS เดิมออกมาแล้ว ทำให้การใช้งานในชีวิตประจำวันไม่มีอุปสรรค อย่างไรก็ตาม แอปเฉพาะทางและแอปอุตสาหกรรมบางส่วนอาจยังอยู่ในระหว่างการปรับใช้ ก่อนอัปเกรดแนะนำให้ตรวจสอบก่อนว่าแอปที่คุณใช้บ่อยรองรับแล้วหรือไม่

\booksubsection{เทคนิคการใช้งาน HarmonyOS ที่ใช้งานได้จริง}

\begin{itemize}
  \item \textbf{เกินระดับสุดท้ายปลาย}: ปัดลงจากมุมขวาบนของหน้าจอเพื่อเปิดศูนย์ควบคุม แตะ "ซูเปอร์เทอร์มินัล" จะเห็นอุปกรณ์อื่นๆ รอบตัว ลากเพื่อเชื่อมต่อ
  \item \textbf{สถานีรับส่งกลาง}: กดค้างที่ข้อความ/รูป/ไฟล์ เลือก "คัดลอกไปยังสถานีรับส่งกลาง" แล้วกดค้างวางในแอปอื่น หรือลากออกมาใช้
  \item \textbf{การ์ดบริการ}: ปัดขึ้นบนไอคอนแอปเพื่อขยายการ์ด กดค้างที่การ์ดเพื่อปรับขนาดหรือเพิ่มลงหน้าจอหลัก
  \item \textbf{หลายหน้าต่างอัจฉริยะ}: ปัดจากขอบซ้ายหรือขวาของหน้าจอเข้าด้านในแล้วค้างไว้ เพื่อเปิดแถบข้าง แตะแอปให้เปิดเป็นหน้าต่างลอย
  \item \textbf{การปกป้องความเป็นส่วนตัว}: เข้า ตั้งค่า > ความเป็นส่วนตัว ดูบันทึกการใช้สิทธิ์ของแอป และตั้งค่าการระบุตำแหน่งแบบเบลอ ฯลฯ
\end{itemize}

\begin{tcbwarn}{ข้อควรระวัง}
ระบบ HarmonyOS อัปเดตต่อเนื่องไม่หยุด แต่ละเวอร์ชันอาจต่างกันบ้าง เนื้อหานี้อธิบายตามคุณสมบัติหลักของ HarmonyOS 4.3 และ 6 ส่วนฟังก์ชันเฉพาะโปรดอ้างอิงจากเวอร์ชันจริงบนโทรศัพท์ของคุณ แนะนำให้อัปเดตระบบเสมอเพื่อรับฟังก์ชันใหม่และแพตช์ความปลอดภัยล่าสุด

\end{tcbwarn}

\begin{tcbpractice}{ลองเลย: สัมผัสฟังก์ชันเด่นของ HarmonyOS}
\begin{enumerate}
  \item เปิดศูนย์ควบคุม หา "ซูเปอร์เทอร์มินัล" ดูว่ารอบๆ มีอุปกรณ์ HarmonyOS อื่นหรือไม่
  \item ลอง "สถานีรับส่งกลาง": คัดลอกข้อความหนึ่งช่วงไปไว้ที่นั่น แล้วลากออกมาในแอปบันทึกช่วยจำ
  \item ปัดขึ้นไอคอนแอปที่ใช้บ่อยสักสองสามตัว ดูว่าการ์ดบริการของมันแสดงอะไร
  \item ไปที่ ตั้งค่า > ความเป็นส่วนตัว ดูว่าล่าสุดแอปไหนใช้สิทธิ์ที่อ่อนไหวบ้าง
\end{enumerate}

\end{tcbpractice}
